\section{Introduction}
\label{sec:intro}

Point cloud registration estimates the rigid motion between two scans. Correspondence-based methods match local descriptors and then estimate the motion robustly from the putative matches. At low overlap, most of these matches are wrong: with FPFH~\cite{rusu2009fpfh} on 3DLoMatch~\cite{huang2021predator}, 888 of the 1781 pairs have fewer than 1\% correct matches (\cref{fig:dilute}). Robust estimators such as RANSAC~\cite{fischler1981ransac}, SC$^2$-PCR~\cite{chen2022sc2}, MAC~\cite{zhang2023mac} and TurboReg~\cite{yan2025turboreg} generate many pose hypotheses in this setting, but each returns only one of them.

A wrong output does not mean that every hypothesis is wrong. The six frontends we evaluate register 35.1 to 40.4\% of the 3DLoMatch+FPFH pairs natively (\cref{tab:main}). After up to 32 of their hypotheses are refined on the observed surfaces, a correct pose is present for 54.2 to 59.8\% of the pairs (\cref{tab:joint}). Each frontend thus loses 16.1 to 20.8 percentage points between its refined pool and its native output.

We therefore treat robust registration as two linked problems. \emph{Recovery} generates and preserves hypotheses that can be refined to the correct pose. \emph{Conversion} spends a limited refinement budget on these hypotheses and selects the final pose from the refined results. For conversion, we build a training-free backend that can follow any frontend. Multi-stream finishing allocates at most 32 point-to-plane ICP refinements across distinct candidates using correspondence support, surface agreement and geometric information. Release then selects from the refined pool by combining geometric support with the concentration of the residuals.

For recovery, we organize the generator after the unit operations that recover a scarce product in chemical separation~\cite{mccabe2005,mullin2001} (\cref{fig:flowsheet}). Each operation addresses one difficulty of generation from contaminated matches. Motif nucleation seeds candidate phases from compatible pairs or triples of matches. A finite site capacity, borrowed from adsorption, penalizes matches that many seeds already use, so that the seeds spread over several clusters. Growth that retains each phase's current composition, as in total reflux, enriches the phase with compatible matches. Together these operations form the NGRR-LS frontend. The analogy guides which operations the process contains, and we test each of them by removing or replacing it.

With the same settings for every frontend, the backend improves all six frontends by 10.8 to 13.7 percentage points on 3DLoMatch+FPFH, and its finishing stage takes 0.09 to 0.11\,s per pair on 16 CPU threads. With NGRR-LS, recall rises from 40.4\% for the native output to 49.6\% after finishing and 53.1\% after release, above the published results collected in \cref{tab:literature}. The compiled pipeline has a median time of 0.39\,s per pair, against 0.92\,s for SC$^2$-PCR with the same backend. Capacity raises the refined-pool coverage from 56.4\% to 59.2\% while native recall changes only from 40.3\% to 40.4\%, and the backend converts this coverage into a 2.0-point gain in final recall. With the same number of ICP calls on every frontend--input cell, finishing registers more pairs than a matched top-$K$ allocation (+693/$-$221 cells).

\paragraph{Contributions.} We make three contributions. (1) A training-free finishing and release backend that allocates a fixed surface-refinement budget across candidates and selects from the refined pool with complementary evidence. It improves six registration frontends at a finishing cost of 0.09 to 0.11\,s per pair. (2) The NGRR-LS frontend, which combines motif nucleation, soft site capacity and retained growth so that the pool contains a correct, refinable candidate more often. (3) An evaluation that separates candidate coverage from final selection, with stage diagnostics, interventions on each operation, controls at a matched ICP budget, nine real-data cohorts and CPU timing.
